% !Mode:: "TeX:UTF-8"
% !TEX program  = xelatex

\documentclass[bwprint]{report}

% ---- 流程图与代码排版所需宏包 ----
\usepackage{tikz}
\usetikzlibrary{shapes.geometric, arrows.meta, positioning, calc}
\usepackage[section]{placeins}
\usepackage{flafter}

% 本报告使用网络空间安全课程设计内容，禁用模板中自带的“视觉计算课程设计任务书”页。
\renewcommand{\makeothertitle}{\clearpage}

% 代码清单样式（Python）
\definecolor{codekw}{RGB}{0,0,180}
\definecolor{codecmt}{RGB}{0,128,0}
\definecolor{codestr}{RGB}{163,21,21}
\definecolor{codebg}{RGB}{248,248,248}
\lstset{
  basicstyle=\ttfamily\small,
  keywordstyle=\color{codekw}\bfseries,
  commentstyle=\color{codecmt},
  stringstyle=\color{codestr},
  numbers=left, numberstyle=\tiny\color{gray}, numbersep=6pt,
  backgroundcolor=\color{codebg},
  frame=single, rulecolor=\color{gray!40},
  breaklines=true, showstringspaces=false,
  tabsize=4, captionpos=b, language=Python,
  xleftmargin=1.2em,
}

% 流程图节点样式
\tikzset{
  fcbase/.style={draw, align=center, font=\small, minimum height=7mm, inner sep=3pt},
  startstop/.style={fcbase, rounded corners=10pt, fill=blue!12},
  process/.style={fcbase, rectangle, fill=gray!8},
  io/.style={fcbase, trapezium, trapezium left angle=70, trapezium right angle=110, fill=green!10},
  decision/.style={fcbase, diamond, aspect=2.2, fill=orange!18, inner sep=1pt},
  fcarrow/.style={-{Stealth[length=2mm]}, thick},
}

% 调整目录排版，给中文章节号和标题留出更宽的空间，避免目录显得拥挤。
\setlength{\cftbeforesecskip}{10pt}
\setlength{\cftbeforesubsecskip}{5pt}
\setlength{\cftsecnumwidth}{4.8em}
\renewcommand{\cftsecaftersnumb}{\hskip1.0em}
\setlength{\cftsubsecindent}{5.8em}
\setlength{\cftsubsecnumwidth}{4.2em}

\title{Web 登录系统安全加固与审计日志设计}
\stunum{2023112621}
\stuname{程泽凯}
\stuclass{计算机类 2023 级}
\supervisor{郑丹阳 \hspace{1.0em} 老师}
\dateinput{\today}

\begin{document}

\maketitle

% 目录使用 Roman 页脚
\pagestyle{plain}
\pagenumbering{Roman}
\tableofcontents
\newpage
\pagestyle{plain}

% 正文页脚
\setcounter{page}{1}
\pagenumbering{arabic}


\section{课程设计题目}
\textbf{题目：Web 登录系统安全加固与审计日志设计。}

本课程设计由本人程泽凯（学号 2023112621）独立完成。围绕一个本地运行的小型 Web 登录系统，本人完成了需求分析、技术选型、数据库与安全规则设计、Python 编码、模拟数据构造、测试用例编写、流程图绘制、实验验证以及 \LaTeX{} 报告排版。项目重点不是简单实现登录页面，而是在可运行的登录流程中加入口令安全存储、登录二次验证模拟、失败限制、账号锁定、审计留痕和异常识别机制，并通过本人编写的测试脚本验证设计效果。

\textbf{合规声明：} 本课程设计全部在本机（Windows 11 家庭中文版 + Python 3.13）离线环境与本地浏览器中完成。数据库、登录日志、测试账号与异常登录记录均为本人自行生成的模拟数据，\textbf{未对任何真实第三方系统、公网服务器或校园网设备进行扫描、攻击或测试，未造成任何影响。}


\section{项目背景}
登录系统是 Web 应用的安全入口，也是攻击者的首要目标。一旦登录环节被攻破，攻击者即可获得合法身份并进一步访问受保护资源。常见的登录安全威胁包括：

\begin{itemize}
    \item \textbf{口令明文或弱哈希存储：} 数据库一旦泄露，口令直接暴露，并可能被用于撞库攻击；
    \item \textbf{暴力破解与撞库：} 攻击者借助自动化脚本对账号进行高频口令尝试；
    \item \textbf{异常登录缺乏感知：} 异地登录、非工作时间登录、高频请求等可疑行为得不到及时识别；
    \item \textbf{缺乏审计能力：} 登录行为无日志记录，安全事件发生后无法追溯定位。
\end{itemize}

针对上述问题，OWASP 在其认证安全指南中明确建议对口令进行加盐慢哈希存储、对失败登录进行限制并记录审计日志\cite{owasp2021authn}；NIST SP 800-63B 数字身份指南同样强调了口令存储与认证生命周期管理的安全要求\cite{nist2020sp80063b}。本课程设计正是围绕这些安全实践，构建一个具备\textbf{安全加固与审计能力}的小型登录系统，做到“有场景、有数据、有流程、有规则、有实现、有测试、有结果、有分析”。

\section{实验环境}
本系统的开发与测试环境如表~\ref{tab:env} 所示。选型上采用 Python Flask + SQLite，原因在于二者轻量、零外部服务依赖，便于本地一键运行与答辩现场演示；同时不引入 ORM，所有 SQL 透明可审计，便于讲解与安全分析。

\begin{table}[H]
    \caption{实验环境配置}
    \label{tab:env}
    \centering
    \begin{tabular}{ll}
        \toprule[1.5pt]
        类别 & 配置 \\
        \midrule[1pt]
        操作系统 & Windows 11 家庭中文版（10.0.26200） \\
        开发语言 & Python 3.13.5 \\
        Web 框架 & Flask 3.1.0 / Werkzeug 3.1.3 \\
        数据库 & SQLite 3（Python 标准库 \texttt{sqlite3}） \\
        口令哈希 & PBKDF2-SHA256，26 万次迭代，自带随机盐 \\
        运行与测试 & 本地浏览器（127.0.0.1:5000）+ 自动化测试脚本 \\
        \bottomrule[1.5pt]
    \end{tabular}
\end{table}

图~\ref{fig:env} 为本人在本机 PowerShell 中执行 \texttt{python --version} 与 \texttt{pip show flask} 后截取的实验环境信息。

\begin{figure}[H]
    \centering
    \includegraphics[width=0.92\textwidth]{2026-06-02_16-57-31.png}
    \caption{本人本机 Python 与 Flask 实验环境}
    \label{fig:env}
\end{figure}

\section{需求分析}
\subsection{功能需求}
任务书要求“必须实现的功能”共 7 项，本系统逐条实现并定位到对应代码模块，如表~\ref{tab:func} 所示。

\begin{table}[H]
    \caption{必做功能与实现位置对照}
    \label{tab:func}
    \centering
    \begin{tabular}{clc}
        \toprule[1.5pt]
        序号 & 必须实现的功能 & 实现位置 \\
        \midrule[1pt]
        1 & 用户注册 & \texttt{app.py} / \texttt{security.register\_user} \\
        2 & 用户登录 & \texttt{app.py} / \texttt{security.authenticate} \\
        3 & 密码哈希存储 & \texttt{security.hash\_password} \\
        4 & 登录失败计数 & \texttt{security.count\_recent\_failures} \\
        5 & 账号临时锁定 & \texttt{security} 锁定逻辑 + \texttt{locked\_until} \\
        6 & 登录日志记录 & \texttt{db.add\_login\_log} + \texttt{login\_logs} 表 \\
        7 & 异常登录提醒 & \texttt{audit.detect\_realtime} \\
        \bottomrule[1.5pt]
    \end{tabular}
\end{table}

\subsection{安全需求}
\begin{itemize}
    \item 口令绝不明文存储，采用加盐慢哈希；
    \item 用户名不存在与口令错误返回统一提示，防止用户名枚举；
    \item 所有数据库操作使用参数化查询（占位符 \texttt{?}），防止 SQL 注入；
    \item 失败登录达阈值后临时锁定账号，抵抗暴力破解；
    \item 登录页加入固定码 \texttt{12345} 的二次验证模拟，只有用户名、口令与验证码均通过后才进入认证成功流程；
    \item 审计可重复执行并输出结构化结果，便于事后追溯。
\end{itemize}

\subsection{异常规则需求与个性化阈值设计}
任务书给出四类异常的规则示例，本人结合自身上机作息与实验场景，将其细化为可量化、可现场修改的阈值，如表~\ref{tab:rules} 所示。所有阈值集中存放于 \texttt{config.py}，答辩时可现场修改任一数值并重新运行验证。

\begin{table}[H]
    \caption{四类异常的个性化判断规则（设计者：程泽凯）}
    \label{tab:rules}
    \centering
    \begin{tabular}{lll}
        \toprule[1.5pt]
        异常类型 & 任务书规则示例 & 本人设计的判断规则 \\
        \midrule[1pt]
        暴力尝试 & 5 分钟内失败超过 5 次 & 5 分钟窗口内失败 $\geq 5$ 次 $\rightarrow$ 锁定 10 分钟 \\
        异地登录 & 同账号短时不同 IP & 10 分钟内同账号出现 $\geq 2$ 个不同 IP \\
        非工作时间 & 晚 23:00 后登录 & 23:00--06:00 之间成功登录 \\
        高频登录 & 1 分钟内超过阈值 & 60 秒内请求 $\geq 10$ 次 \\
        \bottomrule[1.5pt]
    \end{tabular}
\end{table}

\section{技术方案}
\subsection{系统总体架构}
系统采用分层设计，由配置层、数据访问层、安全核心层、审计层与 Web 表现层构成。各模块职责单一、低耦合，所有安全阈值与模拟二次验证码集中于配置层。整体调用关系如下：浏览器经 HTTP 访问 Flask 路由层（\texttt{app.py}），路由层先校验登录页输入的二次验证模拟码 \texttt{12345}，再调用安全核心（\texttt{security.py}）完成认证，调用审计模块（\texttt{audit.py}）完成异常检测，二者均通过数据访问层（\texttt{db.py}）读写 SQLite 数据库。

图~\ref{fig:tree} 展示了本人本地项目目录。代码按配置、数据访问、安全认证、异常审计、页面模板与自动化测试拆分，便于逐模块验证。

\begin{figure}[H]
    \centering
    \includegraphics[height=0.60\textheight]{2026-06-02_16-58-25.png}
    \caption{本人本地项目目录结构}
    \label{fig:tree}
\end{figure}

\subsection{密码哈希方案}
系统采用 Werkzeug 提供的 \texttt{pbkdf2:sha256}（26 万次迭代）对口令进行加盐慢哈希存储，符合 PKCS \#5 / RFC 8018 对基于口令的密钥派生函数的规范\cite{rfc8018}。其安全意义在于三点：（1）\textbf{加盐}——每个口令使用独立随机盐，相同口令的哈希值各不相同，可抵抗彩虹表攻击；（2）\textbf{慢哈希}——26 万次迭代显著抬高离线爆破成本；（3）\textbf{恒定时间比较}——校验时使用恒定时间比较，防止时序侧信道泄露。核心实现如代码清单~\ref{lst:hash} 所示。

\begin{lstlisting}[caption={口令哈希与校验（security.py 节选）}, label={lst:hash}]
from werkzeug.security import generate_password_hash, check_password_hash

# Store: salted slow hash, never plaintext
def hash_password(plain):
    return generate_password_hash(plain, method="pbkdf2:sha256:260000")

# Verify: constant-time comparison
def verify_password(stored_hash, plain):
    return check_password_hash(stored_hash, plain)
\end{lstlisting}

\subsection{数据库设计}
数据库包含三张表：用户表 \texttt{users}、登录日志表 \texttt{login\_logs} 与审计告警表 \texttt{audit\_alerts}，并为日志的高频查询字段建立联合索引，如代码清单~\ref{lst:schema} 所示。其中 \texttt{login\_logs.reason} 取值为 \texttt{ok / wrong\_password / user\_not\_found / locked}，是后续审计判断的关键字段。

\begin{lstlisting}[caption={数据库表结构（db.py 节选）}, label={lst:schema}, language=SQL]
CREATE TABLE users (
  id INTEGER PRIMARY KEY AUTOINCREMENT,
  username TEXT NOT NULL UNIQUE,
  password_hash TEXT NOT NULL,
  created_at TEXT NOT NULL,
  failed_count INTEGER NOT NULL DEFAULT 0,
  locked_until TEXT,           -- lock deadline, NULL = not locked
  last_login_at TEXT
);
CREATE TABLE login_logs (
  id INTEGER PRIMARY KEY AUTOINCREMENT,
  username TEXT NOT NULL, ip TEXT NOT NULL, user_agent TEXT,
  success INTEGER NOT NULL,    -- 1 success / 0 failure
  reason TEXT NOT NULL, created_at TEXT NOT NULL
);
CREATE TABLE audit_alerts (
  id INTEGER PRIMARY KEY AUTOINCREMENT,
  username TEXT NOT NULL, alert_type TEXT NOT NULL,
  detail TEXT NOT NULL, created_at TEXT NOT NULL
);
CREATE INDEX idx_logs_user_time ON login_logs(username, created_at);
\end{lstlisting}

\subsection{登录失败限制、二次验证与异常检测策略}
失败计数采用\textbf{基于日志的时间窗口统计}而非简单计数器累加，因此“5 分钟内失败 5 次”这一时间语义是精确的。当窗口内失败次数达到阈值，即写入锁定截止时间并生成暴力尝试告警。为满足登录加固中的验证码或二次验证要求，本人另外在登录页加入固定码 \texttt{12345} 的二次验证模拟：验证码错误时，系统不进入后续口令校验与会话写入，直接提示“二次验证失败”；验证码正确时，才继续调用 \texttt{security.authenticate()}。异常检测分为两个阶段：登录成功瞬间运行轻量的\textbf{实时检测}（异地 / 非工作时间 / 高频），命中即弹出提醒；离线运行\textbf{全量审计}，按时间回放全部日志、对四类规则逐条匹配并去重，输出结构化报告。

\section{系统总体流程图}
为使系统结构与本人编写的代码能够相互对应，本节给出系统总体处理流程。浏览器请求首先进入 \texttt{app.py} 路由层，再根据用户选择进入注册、登录或安全审计功能；登录认证与异常审计的细化步骤将在后续章节展开。

图~\ref{fig:overall} 给出系统的总体处理流程，对应 \texttt{app.py} 的路由分发。

\begin{figure}[H]
    \centering
    \begin{tikzpicture}[node distance=9mm]
        \node[startstop] (a) {用户访问系统};
        \node[decision, below=of a] (b) {已登录?};
        \node[process, right=25mm of b] (c) {进入个人中心\\/dashboard};
        \node[decision, below=of b] (d) {选择操作};
        \node[io, below left=12mm and 8mm of d] (e) {注册：输入信息};
        \node[io, below right=12mm and 8mm of d] (f) {登录：输入信息\\含验证码 12345};
        \node[process, below=18mm of d] (g) {输入校验与认证\\注册时哈希存储\\登录时校验 12345\\再执行图\ref{fig:auth}};
        \node[startstop, below=of g] (h) {进入个人中心或安全审计\\audit（见图\ref{fig:audit}）};

        \draw[fcarrow] (a) -- (b);
        \draw[fcarrow] (b) -- node[above]{是} (c);
        \draw[fcarrow] (b) -- node[right]{否} (d);
        \draw[fcarrow] (d) -| node[near start, above]{注册} (e);
        \draw[fcarrow] (d) -| node[near start, above]{登录} (f);
        \draw[fcarrow] (e) |- (g);
        \draw[fcarrow] (f) |- (g);
        \draw[fcarrow] (g) -- (h);
    \end{tikzpicture}
    \caption{系统总体流程图}
    \label{fig:overall}
\end{figure}

\section{核心模块流程图}
本节给出两个关键安全模块的流程：用户登录认证与登录失败限制。两张流程图均由本人根据实际代码分支绘制，可用于答辩时逐项说明“流程图、代码与安全规则”的对应关系。

\subsection{用户登录认证流程}
图~\ref{fig:auth} 对应 \texttt{security.authenticate()} 函数，完整体现“查询用户 $\rightarrow$ 判断锁定 $\rightarrow$ 判断用户是否存在 $\rightarrow$ 校验口令 $\rightarrow$ 成功后实时异常检测”的判断分支与异常路径。

\begin{figure}[H]
    \centering
    \begin{tikzpicture}[node distance=8mm]
        \node[startstop] (s) {输入用户名+口令};
        \node[process, below=of s] (q) {查询用户 db.get\_user};
        \node[decision, below=of q] (lock) {账号锁定?};
        \node[process, right=24mm of lock] (lockr) {记 locked 日志\\$\rightarrow$ 拒绝};
        \node[decision, below=10mm of lock] (exist) {用户存在?};
        \node[process, right=24mm of exist] (existr) {记 user\_not\_found\\$\rightarrow$ 统一提示};
        \node[decision, below=10mm of exist] (pwd) {口令正确?};
        \node[process, right=24mm of pwd] (pwdr) {记 wrong\_password\\$\rightarrow$ 失败限制(图\ref{fig:lock})};
        \node[process, below=10mm of pwd] (ok) {记 ok 日志\\清零失败计数};
        \node[process, below=of ok] (rt) {实时异常检测\\detect\_realtime};
        \node[decision, below=of rt] (al) {命中异常?};
        \node[process, right=24mm of al] (alr) {弹出异常\\登录提醒};
        \node[startstop, below=12mm of al] (end) {登录成功，写入会话};

        \draw[fcarrow] (s)--(q);
        \draw[fcarrow] (q)--(lock);
        \draw[fcarrow] (lock)--node[above]{是}(lockr);
        \draw[fcarrow] (lock)--node[right]{否}(exist);
        \draw[fcarrow] (exist)--node[above]{否}(existr);
        \draw[fcarrow] (exist)--node[right]{是}(pwd);
        \draw[fcarrow] (pwd)--node[above]{否}(pwdr);
        \draw[fcarrow] (pwd)--node[right]{是}(ok);
        \draw[fcarrow] (ok)--(rt);
        \draw[fcarrow] (rt)--(al);
        \draw[fcarrow] (al)--node[above]{是}(alr);
        \draw[fcarrow] (al)--node[right]{否}(end);
        \draw[fcarrow] (alr) |- (end);
    \end{tikzpicture}
    \caption{用户登录认证流程图}
    \label{fig:auth}
\end{figure}

\subsection{登录失败限制流程}
图~\ref{fig:lock} 对应口令错误分支与 \texttt{count\_recent\_failures()}：失败计数加一、判断阈值、锁定账号、记录告警。

\begin{figure}[H]
    \centering
    \begin{tikzpicture}[node distance=8mm]
        \node[startstop] (s) {口令校验失败};
        \node[process, below=of s] (log) {写入失败日志};
        \node[process, below=of log] (cnt) {统计 5 分钟窗口内\\失败次数};
        \node[process, below=of cnt] (upd) {更新 failed\_count};
        \node[decision, below=of upd] (dec) {失败次数 $\geq 5$?};
        \node[startstop, right=26mm of dec] (no) {提示：还可\\尝试 N 次};
        \node[process, below=10mm of dec] (calc) {计算锁定截止\\now + 10 分钟};
        \node[process, below=of calc] (set) {写入 locked\_until};
        \node[process, below=of set] (alert) {生成暴力尝试告警};
        \node[startstop, below=of alert] (end) {提示：账号已锁定};

        \draw[fcarrow] (s)--(log);
        \draw[fcarrow] (log)--(cnt);
        \draw[fcarrow] (cnt)--(upd);
        \draw[fcarrow] (upd)--(dec);
        \draw[fcarrow] (dec)--node[above]{否}(no);
        \draw[fcarrow] (dec)--node[right]{是}(calc);
        \draw[fcarrow] (calc)--(set);
        \draw[fcarrow] (set)--(alert);
        \draw[fcarrow] (alert)--(end);
    \end{tikzpicture}
    \caption{登录失败限制流程图}
    \label{fig:lock}
\end{figure}

\section{测试与处置流程图}
本节给出异常登录审计流程与测试处置流程。前者对应本人实现的 \texttt{audit.py}，后者对应本人编写的 \texttt{run\_tests.py} 与异常处置策略。通过将“检测”和“处置”连成闭环，可以更清楚地体现系统的实际安全用途。

\subsection{异常登录审计流程}
图~\ref{fig:audit} 对应 \texttt{audit.run\_full\_audit()}：清空旧告警、按时间遍历日志、对四类规则匹配、汇总输出审计结果。

\begin{figure}[H]
    \centering
    \begin{tikzpicture}[node distance=8mm]
        \node[startstop] (s) {开始审计};
        \node[process, below=of s] (clr) {清空旧告警 clear\_alerts};
        \node[process, below=of clr] (read) {按时间读取全部日志};
        \node[decision, below=of read] (loop) {遍历每条日志};
        \node[process, right=26mm of loop, text width=33mm] (rules) {规则匹配：\\失败$\geq$5 暴力尝试\\23-6 点 非工作时间\\10 分钟多 IP 异地\\60 秒$\geq$10 高频};
        \node[process, below=14mm of loop] (sum) {汇总各类告警计数};
        \node[startstop, below=of sum] (out) {输出审计报告\\日志数/成功/失败/告警};

        \draw[fcarrow] (s)--(clr);
        \draw[fcarrow] (clr)--(read);
        \draw[fcarrow] (read)--(loop);
        \draw[fcarrow] (loop)--node[above]{逐条}(rules);
        \draw[fcarrow] (rules) |- (sum);
        \draw[fcarrow] (loop)--node[right]{完成}(sum);
        \draw[fcarrow] (sum)--(out);
    \end{tikzpicture}
    \caption{异常登录审计流程图}
    \label{fig:audit}
\end{figure}

\subsection{自动化测试与异常处置流程}
图~\ref{fig:test-dispose} 对应 \texttt{run\_tests.py} 与系统的异常处置思路。每个测试用例执行前先重建数据库，避免不同用例互相污染；若用例未通过，则返回检查规则、阈值与实现。告警产生后，根据不同异常类型采取锁定、提醒、人工复核或限流等处置方式。

\begin{figure}[H]
    \centering
    \begin{tikzpicture}[node distance=8mm]
        \node[startstop] (s) {运行 run\_tests.py};
        \node[process, below=of s] (fresh) {每个用例重建数据库\\fresh\_db 保证隔离};
        \node[process, below=of fresh] (tc) {执行 TC1--TC9 断言};
        \node[decision, below=of tc] (pass) {断言全部通过?};
        \node[process, right=24mm of pass, text width=31mm] (fix) {定位失败用例\\检查规则、阈值与实现};
        \node[decision, below=12mm of pass] (type) {告警类型?};
        \node[process, below left=13mm and 29mm of type, text width=24mm] (bf) {暴力尝试：\\临时锁定账号};
        \node[process, below left=13mm and 2mm of type, text width=24mm] (ip) {异地登录：\\提醒二次验证};
        \node[process, below right=13mm and 2mm of type, text width=24mm] (off) {非工作时间：\\记录并人工复核};
        \node[process, below right=13mm and 29mm of type, text width=24mm] (freq) {高频登录：\\限流或拉黑 IP};
        \node[startstop, below=38mm of type] (out) {输出处置结论\\与审计报告};

        \draw[fcarrow] (s)--(fresh);
        \draw[fcarrow] (fresh)--(tc);
        \draw[fcarrow] (tc)--(pass);
        \draw[fcarrow] (pass)--node[above]{否}(fix);
        \draw[fcarrow] (fix) |- (fresh);
        \draw[fcarrow] (pass)--node[right]{是}(type);
        \draw[fcarrow] (type) -| node[near start, above]{暴力} (bf);
        \draw[fcarrow] (type) -| node[near start, above]{异地} (ip);
        \draw[fcarrow] (type) -| node[near start, above]{非工时} (off);
        \draw[fcarrow] (type) -| node[near start, above]{高频} (freq);
        \draw[fcarrow] (bf) |- (out);
        \draw[fcarrow] (ip) |- (out);
        \draw[fcarrow] (off) |- (out);
        \draw[fcarrow] (freq) |- (out);
    \end{tikzpicture}
    \caption{自动化测试与异常处置流程图}
    \label{fig:test-dispose}
\end{figure}

\section{系统实现}
\subsection{登录认证核心}
登录认证严格按图~\ref{fig:auth} 的四个分支实现：账号锁定则直接拒绝；用户不存在返回统一提示防枚举；口令错误进入失败限制流程；认证成功后清零计数并运行实时异常检测。核心逻辑如代码清单~\ref{lst:auth} 所示。

\begin{lstlisting}[caption={登录认证主流程（security.authenticate 节选）}, label={lst:auth}]
def authenticate(username, password, ip, user_agent):
    with db.get_conn() as conn:
        user = db.get_user(conn, username)
        # branch 1: account locked -> reject directly
        if user is not None and is_locked(user):
            db.add_login_log(conn, username, ip, user_agent, False, "locked", ts)
            return {"success": False, "reason": "locked", ...}
        # branch 2: user not found -> unified message (anti-enumeration)
        if user is None:
            db.add_login_log(conn, username, ip, user_agent, False, "user_not_found", ts)
            return {"success": False, "reason": "user_not_found", ...}
        # branch 3: wrong password -> count + maybe lock
        if not verify_password(user["password_hash"], password):
            db.add_login_log(conn, username, ip, user_agent, False, "wrong_password", ts)
            failures = count_recent_failures(conn, username)
            if failures >= MAX_FAILED_ATTEMPTS:
                db.set_lock(conn, username, lock_until)   # lock 10 min
                db.add_alert(conn, username, "brute_force", ...)
            return {"success": False, "reason": "wrong_password", ...}
        # branch 4: success -> reset + realtime anomaly detection
        db.add_login_log(conn, username, ip, user_agent, True, "ok", ts)
        db.reset_user_security(conn, username, ts)
        alerts = audit.detect_realtime(conn, username, ip, now())
    return {"success": True, "reason": "ok", "alerts": alerts}
\end{lstlisting}

\subsection{异常检测实现}
非工作时间判断需处理跨午夜区间（23:00 至次日 06:00），异地登录通过统计窗口内不同 IP 数实现，高频登录通过秒级窗口计数实现，如代码清单~\ref{lst:audit} 所示。

\begin{lstlisting}[caption={异常规则判断（audit.py 节选）}, label={lst:audit}]
def is_off_hours(dt):              # rule 3: cross-midnight interval
    h = dt.hour
    start, end = 23, 6
    return h >= start or h < end

def distinct_ips_in_window(conn, username, end_dt, minutes):  # rule 2
    start = (end_dt - timedelta(minutes=minutes)).isoformat(sep=" ")
    cur = conn.execute("SELECT DISTINCT ip FROM login_logs "
        "WHERE username=? AND created_at>=? AND created_at<=?",
        (username, start, end_dt.isoformat(sep=" ")))
    return [r["ip"] for r in cur.fetchall()]
\end{lstlisting}

\subsection{Web 页面、二次验证与异常提醒}
为便于在答辩现场演示，本人实现了注册页、登录页、个人中心与安全审计页。登录页除用户名和口令外，还加入“二次验证模拟码”输入框，实验固定值为 \texttt{12345}；若输入其他值，\texttt{app.py} 会在调用认证函数前直接拒绝，防止仅凭口令进入系统。登录页还额外提供“模拟来源 IP”输入框，仅用于本地实验中构造异地登录记录。图~\ref{fig:web-auth} 展示了注册页和登录页，图~\ref{fig:multi-ip-ui} 展示了同一账号短时间使用多个模拟 IP 登录后触发的异常提醒。

\begin{lstlisting}[caption={二次验证模拟码校验（app.py 节选）}, label={lst:otp}]
otp_code = request.form.get("otp_code", "").strip()
if otp_code != config.SIMULATED_OTP_CODE:
    flash("二次验证失败：请输入本地实验验证码 12345", "danger")
    return render_template("login.html")
\end{lstlisting}

\begin{figure}[H]
    \centering
    \begin{subfigure}[b]{0.48\textwidth}
        \centering
        \includegraphics[width=\textwidth]{2026-06-02_16-59-08.png}
        \caption{用户注册页面}
    \end{subfigure}
    \hfill
    \begin{subfigure}[b]{0.48\textwidth}
        \centering
        \includegraphics[width=\textwidth]{2026-06-02_17-00-09.png}
        \caption{用户登录页面}
    \end{subfigure}
    \caption{本人实现的 Web 注册与登录页面}
    \label{fig:web-auth}
\end{figure}

\begin{figure}[H]
    \centering
    \includegraphics[width=0.92\textwidth]{2026-06-02_17-02-40.png}
    \caption{本人模拟多个来源 IP 后触发异地登录提醒}
    \label{fig:multi-ip-ui}
\end{figure}

\section{实验数据}
\subsection{测试账号（本人创建）}
本人在 \texttt{seed\_data.py} 中创建 4 个测试账号，口令均经加盐哈希存储，如表~\ref{tab:accounts} 所示。

\begin{table}[H]
    \caption{本人创建的测试账号}
    \label{tab:accounts}
    \centering
    \begin{tabular}{lll}
        \toprule[1.5pt]
        用户名 & 用途 & 口令（已加盐哈希存储） \\
        \midrule[1pt]
        czk2023 & 主账号（程泽凯本人） & Czk@2023112621 \\
        alice99 & 普通用户 & Alice\#pass01 \\
        bob\_dev & 普通用户 & BobDev2024! \\
        victim7 & 演示被暴力破解的受害账号 & Victim\#0001 \\
        \bottomrule[1.5pt]
    \end{tabular}
\end{table}

口令哈希的真实存储样例（库中无任何明文）：
\begin{lstlisting}[numbers=none, language=]
czk2023 -> pbkdf2:sha256:260000$W10OrlyPNRYsjpci$8bfdef18688201d9f...
alice99 -> pbkdf2:sha256:260000$77QdBCxQwoKr8Zoy$01689276b56190d39...
\end{lstlisting}

图~\ref{fig:hash-shot} 为本人从本地 SQLite 数据库查询得到的口令哈希截图。可以看到，数据库只保存带随机盐的 PBKDF2-SHA256 哈希，不保存明文口令。

\begin{figure}[H]
    \centering
    \includegraphics[width=0.98\textwidth]{2026-06-02_17-03-56.png}
    \caption{本人本地数据库中的口令哈希存储结果}
    \label{fig:hash-shot}
\end{figure}

\subsection{模拟异常登录记录（本人设计）}
为覆盖四类异常并保留一组正常对照，本人设计了如表~\ref{tab:scenarios} 所示的模拟数据。

\begin{table}[H]
    \caption{本人设计的模拟登录场景}
    \label{tab:scenarios}
    \centering
    \begin{tabular}{cll}
        \toprule[1.5pt]
        场景 & 账号 & 设计的数据 \\
        \midrule[1pt]
        A 暴力尝试 & victim7 & 3 分钟内连续失败 6 次 \\
        B 异地登录 & alice99 & 4 分钟内来自 10.0.0.5 / 172.16.8.9 / 203.0.113.7 \\
        C 非工作时间 & bob\_dev & 凌晨 02:30 成功登录 \\
        D 高频登录 & czk2023 & 40 秒内尝试 12 次 \\
        E 正常对照 & czk2023 & 白天 14:00 单 IP 低频登录（不应告警） \\
        \bottomrule[1.5pt]
    \end{tabular}
\end{table}

\subsection{审计真实输出}
对上述种子数据运行全量审计，本机真实输出为：日志总数 23 条，其中成功登录 6 条、失败登录 17 条，共生成 10 条告警，分类统计如表~\ref{tab:auditout} 所示。

\begin{table}[H]
    \caption{审计告警分类统计（本机真实结果）}
    \label{tab:auditout}
    \centering
    \begin{tabular}{lc}
        \toprule[1.5pt]
        异常类型 & 命中次数 \\
        \midrule[1pt]
        暴力尝试 & 4 \\
        异地登录 & 2 \\
        非工作时间登录 & 1 \\
        高频登录 & 3 \\
        \midrule[1pt]
        合计 & 10 \\
        \bottomrule[1.5pt]
    \end{tabular}
\end{table}

部分告警明细（真实截取）：
\begin{lstlisting}[numbers=none, language=]
[非工作时间登录] bob_dev  于 2026-05-26 02:30:00 在非工作时段成功登录
[暴力尝试]      victim7  截至 2026-05-26 19:10:56，5 分钟内失败 5 次
[异地登录]      alice99  10 分钟内出现 3 个 IP: 10.0.0.5, 172.16.8.9, 203.0.113.7
[高频登录]      czk2023  60 秒内请求 12 次
\end{lstlisting}

图~\ref{fig:audit-alert-shot} 与图~\ref{fig:audit-log-shot} 为本人本地安全审计页面截图，分别展示告警明细与最近登录日志。页面中可以看到本人实际触发的异地登录、暴力尝试以及多条登录失败记录。

\begin{figure}[H]
    \centering
    \includegraphics[width=0.94\textwidth]{2026-06-02_17-30-35.png}
    \caption{本人本地安全审计页面中的告警明细}
    \label{fig:audit-alert-shot}
\end{figure}

\begin{figure}[H]
    \centering
    \includegraphics[width=0.94\textwidth]{2026-06-02_17-30-48.png}
    \caption{本人本地安全审计页面中的最近登录日志}
    \label{fig:audit-log-shot}
\end{figure}

\section{测试用例}
为验证本人设计的安全规则是否真正生效，我编写了自动化测试脚本 \texttt{run\_tests.py}。每个用例执行前均独立重建数据库以保证隔离，共包含 9 条用例、20 项断言，覆盖正常功能、安全存储、二次验证模拟、失败限制、异常检测与全量审计，如表~\ref{tab:testcase} 所示。

\begin{table}[H]
    \caption{测试用例（$\geq 5$ 条，覆盖不同类型）}
    \label{tab:testcase}
    \centering
    \begin{tabular}{cllc}
        \toprule[1.5pt]
        编号 & 测试用例 & 类型 & 结果 \\
        \midrule[1pt]
        TC1 & 注册+口令哈希（无明文/pbkdf2 前缀/弱口令被拒） & 安全存储 & 通过 \\
        TC2 & 正确口令登录成功 & 正常功能 & 通过 \\
        TC3 & 二次验证模拟码错误被拒，输入 \texttt{12345} 后进入认证流程 & 二次验证 & 通过 \\
        TC4 & 连续失败 5 次触发账号锁定 & 失败限制 & 通过 \\
        TC5 & 锁定期内正确口令仍被拒绝 & 失败限制 & 通过 \\
        TC6 & 异地登录检测（短时多 IP） & 异常检测 & 通过 \\
        TC7 & 非工作时间登录检测 & 异常检测 & 通过 \\
        TC8 & 高频登录检测 & 异常检测 & 通过 \\
        TC9 & 全量审计汇总（覆盖四类异常） & 审计 & 通过 \\
        \bottomrule[1.5pt]
    \end{tabular}
\end{table}

\section{运行结果}
本人于本机重新运行 \texttt{python run\_tests.py}，自动化测试最终输出 \textbf{20/20 项断言通过}。其中，新增的二次验证用例验证了错误验证码会被拒绝，输入本人设置的固定模拟码 \texttt{12345} 后才进入后续登录认证；全量审计用例生成 23 条模拟登录日志，识别出 10 条告警，四类异常均被覆盖。关键终端输出如下：

\begin{lstlisting}[numbers=none, language=]
TC9 Full audit summary
  [PASS] alerts generated: 10
  [PASS] brute force covered
  [PASS] multi-IP login covered
  [PASS] off-hours login covered
  [PASS] high-frequency login covered
  Summary: logs=23, success=6, failed=17, alerts=10

Test result: 20/20 assertions passed
\end{lstlisting}

此外，本人通过 Flask 测试客户端执行冒烟测试，\texttt{/login}、\texttt{/register}、\texttt{/audit}、\texttt{/dashboard} 均返回状态码 200，注册 $\rightarrow$ 登录 $\rightarrow$ 个人中心流程完整贯通。结合测试脚本结果，可以确认系统功能、安全规则与审计输出形成了完整闭环。

图~\ref{fig:lock-shot} 为本人使用错误口令连续登录 5 次后的运行截图。系统提示账号锁定 10 分钟，并同步显示“5 分钟内失败 5 次”的暴力尝试告警。

\begin{figure}[H]
    \centering
    \includegraphics[width=0.92\textwidth]{2026-06-02_17-29-55.png}
    \caption{本人连续输错口令后触发账号临时锁定}
    \label{fig:lock-shot}
\end{figure}

\section{问题与改进}
\subsection{遇到的问题及解决}
\begin{enumerate}
    \item \textbf{时间窗口统计：} 最初使用简单计数器累加失败次数，无法表达“5 分钟内”的时间语义；改为基于登录日志按时间窗口统计后，时间约束变得精确。
    \item \textbf{用户枚举风险：} 用户不存在与口令错误若返回不同提示，会泄露账号是否存在；统一为“用户名或口令错误”后修复该信息泄露问题，符合 OWASP Top 10 中对敏感信息暴露的防护建议\cite{owasp2021top10}。
    \item \textbf{审计可重复性：} 重复审计会重复写入告警；通过审计前 \texttt{clear\_alerts} 加去重集合解决，保证幂等。
    \item \textbf{二次验证模拟：} 为满足任务书中“验证码或二次验证模拟”的要求，本人未接入外部短信或邮箱服务，而是在本地实验环境中设计固定验证码 \texttt{12345}，便于答辩现场复现与自动化测试。
\end{enumerate}

\subsection{改进方向}
\begin{itemize}
    \item 将当前固定码 \texttt{12345} 升级为图形验证码或基于时间的一次性口令（TOTP）二次验证；
    \item 接入 IP 地理库实现真正的异地判定（当前为同账号多 IP 启发式）；
    \item 接入限流中间件，对高频来源 IP 自动限流；
    \item 增加告警推送（邮件 / Webhook）并接入集中式日志分析（如 ELK）。
\end{itemize}

\subsection{答辩可现场验证的参数}
修改 \texttt{config.py} 中任一阈值后重新运行 \texttt{python run\_tests.py} 或刷新 \texttt{/audit} 即可验证效果，例如：将 \texttt{MAX\_FAILED\_ATTEMPTS} 由 5 改为 3，账号锁定将更早触发；将 \texttt{HIGH\_FREQ\_THRESHOLD} 由 10 改为 5，高频告警数量随之上升。

\section{总结}
本课程设计由本人独立完成，最终实现了一个具备安全加固与审计能力的 Web 登录系统，覆盖任务书要求的全部 7 项必做功能与 4 类异常规则，做到了“有实现、有数据、有流程图、有验证”。系统以 PBKDF2 加盐慢哈希存储口令，加入固定码 \texttt{12345} 的二次验证模拟，对失败登录进行基于时间窗口的精确计数并自动锁定账号，记录完整的登录审计日志，并通过实时检测与离线全量审计识别暴力尝试、异地登录、非工作时间登录与高频登录四类异常。

在完成过程中，本人不仅实现了页面与后端功能，还自行设计了集中可调的安全阈值、构造了覆盖正常与异常行为的模拟数据、绘制了五张与代码分支对应的流程图，并编写自动化测试脚本完成验证。最终 20/20 项断言全部通过。通过本次课程设计，本人进一步理解了 Web 认证安全、口令存储、登录审计、异常行为检测以及测试验证之间的关系，也认识到安全系统需要将预防、检测、记录和处置结合起来，而不能只停留在单一功能实现上。

\section{附录：代码、配置、日志、截图}
\subsection{项目代码索引}
本课程设计代码均由本人在本地完成，主要文件及作用如表~\ref{tab:appendix-files} 所示。

\begin{table}[H]
    \caption{项目代码与材料索引}
    \label{tab:appendix-files}
    \centering
    \begin{tabular}{ll}
        \toprule[1.5pt]
        文件 & 作用 \\
        \midrule[1pt]
        \texttt{config.py} & 全局配置与本人设计的安全阈值 \\
        \texttt{db.py} & SQLite 建表与参数化数据访问 \\
        \texttt{security.py} & 注册、哈希、失败限制、账号锁定与认证 \\
        \texttt{audit.py} & 实时异常检测与离线全量审计 \\
        \texttt{app.py} & Flask 主程序、路由与页面渲染 \\
        \texttt{seed\_data.py} & 本人设计的测试账号与模拟异常日志 \\
        \texttt{run\_tests.py} & 9 条自动化测试用例、20 项断言 \\
        \texttt{docs/截图清单.md} & 本人实验截图的获取步骤与放置建议 \\
        \bottomrule[1.5pt]
    \end{tabular}
\end{table}

\subsection{关键配置}
以下配置集中体现了本人针对课程设计场景设置的安全阈值。答辩时可修改任一参数并重新运行测试。

\begin{lstlisting}[caption={本人设计的安全阈值（config.py 节选）}, label={lst:config}]
BRUTE_FORCE_WINDOW_MINUTES = 5
MAX_FAILED_ATTEMPTS = 5
ACCOUNT_LOCK_MINUTES = 10
SIMULATED_OTP_CODE = "12345"
MULTI_IP_WINDOW_MINUTES = 10
MULTI_IP_THRESHOLD = 2
OFF_HOURS_START = 23
OFF_HOURS_END = 6
HIGH_FREQ_WINDOW_SECONDS = 60
HIGH_FREQ_THRESHOLD = 10
\end{lstlisting}

\bibliographystyle{plain}
\bibliography{mybib}

\end{document}

